\scanpage{104}
Pour \unsure{chaque} \deleted{strate}\inserted{comp. conn.} $T_\alpha$ de $X'_T$, considérons
\[
\widehat\varphi_\alpha:\widehat\pi_1(X'_{\overline K},T_\alpha)\longrightarrow\widehat\pi_1(X,T)
\]
(c'est un hom injectif si $X'$ rev. étale de $X$, et alors les \deleted{\unsure{différentes im.}} images des $\varphi_\alpha$ sont les conjugués de l'image de l'un d'eux). Donc \inserted{si \deleted{\unsure{on peut}} la condition sur $x'_\alpha,y'_\alpha$ est vraie $\forall\alpha$,} la formule implique ``la'' condition \unsure{intrinsèque}
\[
g^T_{y,x}(u)\in\bigcap_{\alpha\in\pi_0(X'_T)}\operatorname{Im}\widehat\varphi_\alpha
\qquad\forall u\in\Gamma=\Gamma_{\overline K/K}.
\]
\margin{Si $X'$ est un rev. étale de $X$, \unsure{l'ens} \inserted{$E=\pi_0(X'_T)$} \unsure{discret} \deleted{est} \unsure{d'une} ``fibre'' qui décrit $X'$ en termes de $\pi=\pi_1(X,T)$, via la ``\unsure{fibration}'' \unsure{$E$} définie par $T\to X$ ($T$ 1-connexe) -- une action de $\pi$ sur $E$, et le 2° mb de (55), est le noyau $\unsure{\mathcal I|\mathfrak I}$ sous-groupe \inserted{dist.} de $\pi$ en $\pi\longrightarrow\unsure{\mathfrak S_E}$.}

C'est là une condition arithmétique, nécessaire pour que le \unsure{revê.\textsuperscript{t}} fini \inserted{étale} topologique $X'$ de $X$ (\unsure{quitte se borner} à ce cas) puisse \deleted{\unsure{se représenter}} être muni d'une \deleted{$\unsure{K}$} structure de $K$-schéma, pour laquelle $f:X'\to X$ soit \unsure{rat.|rel.}$ /K$, et telle que \deleted{$\unsure{x,y}$}\inserted{$\forall\alpha\in\pi_0(X'_T)$,} $\forall u\in\Gamma$, $u(x'_\alpha)$ et $u(y'_\alpha)$ soient $T$-adjacents. (Ce qui est une condition nettement plus faible que celle que $x',y'\in X'_K$ \inserted{i.e. $x',y'$} fixes sous $\Gamma$)\ldots
\mspage{103}
